% bab6.tex - Kesimpulan dan Saran.
% Semua angka kinerja memakai makro dari tabel/angka.tex (dibuat skrip/tabel_ke_tex.py).
\chapter{KESIMPULAN DAN SARAN}

\section{Kesimpulan}
Berdasarkan hasil pada Bab~V, tiga rumusan masalah pada Bab~I dijawab sebagai berikut.
\begin{enumerate}
    \item \textbf{Peta awal dari citra beresolusi tinggi.} Proyek ini menghasilkan aplikasi pemetaan tutupan lahan yang utuh. Antarmuka web memanggil program inti \kode{landcover}, lalu menampilkan peta kelompok, pelabelan kelas, serta persentase dan luas setiap kelas. Pada potongan Sentinel-2 T49MDM 2048$\times$2048 piksel, aplikasi menghasilkan peta dengan $K=8$ dan tabel luas untuk \SValidHa~ha lahan valid (Subbab~\ref{sec:hasilsentinel}). Pada citra sintetis terbesar yang diukur (7680$\times$4320, sekitar 33 juta piksel, $K=8$), mode Optimalisasi menurunkan Total dari \LoTotalDelapan~ms menjadi \LdLdTotalDelapan~ms. Total ini tidak memuat pemanasan konteks CUDA (\WarmMin--\WarmMax~ms pada sesi pengukuran), dan waktu Proses untuk ukuran ini tidak terekam. Satu \emph{tile} penuh (sekitar 120 juta piksel) tidak diuji karena datanya tidak tersedia. Karena itu, apakah satu \emph{tile} cukup cepat untuk dicoba secara interaktif dengan beberapa nilai $K$ belum dapat disimpulkan.
    \item \textbf{\emph{Bottleneck} dan penempatan tahap.} \emph{Profiling} L0 menunjukkan bahwa K-Means adalah \emph{bottleneck} utama, dengan porsi \PorsiKMLoMin--\PorsiKMLoMax{} dari Total di Mesin B dan \PorsiKMAMin--\PorsiKMAMax{} di Mesin A (Subbab~\ref{sec:profil}). Penempatan akhir setiap tahap ditunjukkan pada Tabel~\ref{tab:penempatan}: K-Means berjalan di CUDA bila $N\times K\geq70$ juta dan di OpenMP bila kurang dari itu; \emph{counts}, pewarnaan, dan tulis PNG memakai OpenMP; baca PNG dan inisialisasi tetap serial. Aturan $N\times K$ memilih konfigurasi tercepat menurut Proses pada \AmbBaruBenarProses{} dari \AmbJumlah{} konfigurasi uji, sedangkan aturan lama hanya pada \AmbLamaBenarProses{} konfigurasi (Subbab~\ref{sec:hasilambang}). Sesudah optimasi, \emph{bottleneck} berpindah ke baca PNG serial, yang mengambil \PorsiBacaLdLdDelapan{} dari Total L1+L2 pada 7680$\times$4320. Mempercepat K-Means lebih jauh hanya dapat memperbaiki Total paling banyak \AmdBatasLdLd{} lagi.
    \item \textbf{Peningkatan kinerja, sumbangan setiap langkah, dan kesamaan hasil.} Pada 7680$\times$4320 dengan $K=8$, \emph{speedup} Total terhadap L0 adalah \LdSTotalDelapan{} untuk L1 (CUDA), \LdbSTotalDelapan{} untuk L2 (OpenMP), dan \LdLdSTotalDelapan{} untuk L1+L2 (Subbab~\ref{sec:bertahap}). Untuk K-Means saja pada $K=8$, CUDA mencapai \HDuaCudaMin--\HDuaCudaMax{} dan OpenMP 8 \emph{thread} \HDuaOmpMin--\HDuaOmpMax{} di semua ukuran yang diukur. L1 saja tertahan di dekat batas Amdahl \AmdBatas, sedangkan L1+L2 melampauinya karena L2 juga memangkas tahap di luar K-Means. Bila pemanasan CUDA ikut dihitung (Proses, 4000$\times$3000), \emph{speedup} konfigurasi tercepat berkisar dari \SProsesMin{} pada $K=2$ sampai \SProsesMax{} pada $K=16$, sehingga keuntungan GPU pada $K$ kecil sebagian besar hilang. Seluruh pemeriksaan pada Tabel~\ref{tab:verifikasi} menunjukkan label, \emph{centroid}, dan \emph{counts} yang identik antarmode, termasuk \kode{labels.u8} Sentinel-2 pada mode cpu, omp, dan cuda. Dengan demikian, persentase dan luas pada kedua mode sama, dan selisih waktu murni berasal dari optimasi.
\end{enumerate}

Terhadap hipotesis awal (Tabel~\ref{tab:hipotesis}): H1 dan H4 terbukti. H2a tercapai pada ukuran terdekat (4000$\times$3000), karena 4096$^2$ tidak diukur. H2b tercapai pada $K=8$, tetapi tidak pada $K\leq4$. H3 tidak dapat disimpulkan untuk \emph{tile} penuh; pada 12 juta piksel dengan $K=8$, target tercapai menurut Total, tetapi tidak menurut Proses.

\section{Saran}
\begin{enumerate}
    \item \textbf{L3 di dalam \emph{pipeline}.} Varian \emph{streams} perlu dipanggil dari \kode{landcover} dan digabung dengan pemrosesan banyak wilayah, sehingga baca citra wilayah berikutnya di CPU berjalan bersamaan dengan K-Means di GPU. Sebelum itu, penyebab varian \emph{streams} lebih lambat daripada T4 di tingkat \emph{kernel} perlu diselidiki, misalnya dengan mengalokasikan memori \emph{pinned} sekali lalu memakainya ulang.
    \item \textbf{Baca PNG paralel.} Baca PNG adalah \emph{bottleneck} berikutnya. Masukan berupa folder potongan dengan \emph{manifest} sudah dibaca paralel per potongan oleh \kode{landcover} (Tabel~\ref{tab:penempatan}), tetapi belum dipakai pada \emph{benchmark}. Manfaatnya perlu diukur, dan persiapan data perlu langsung menghasilkan citra dalam bentuk potongan.
    \item \textbf{Proses CUDA yang menetap.} Bila aplikasi menjaga satu proses CUDA tetap hidup, pemanasan konteks hanya dibayar sekali. Ambang mode otomatis kemudian perlu dihitung ulang dengan Total, bukan Proses.
    \item \textbf{\emph{Profiling} dengan alat khusus.} Nsight Systems dan Nsight Compute dapat menjelaskan rincian \emph{kernel} dan transfer CUDA, sedangkan \emph{profiler} CPU dapat menjelaskan turunnya efisiensi OpenMP di atas 8 \emph{thread}. Beberapa sebab di Bab~V masih berupa dugaan.
    \item \textbf{Uji \emph{tile} penuh.} Pengukuran sebaiknya diulang pada \emph{tile} Sentinel-2 asli 10980$\times$10980 dengan minimal 10 ulangan sesuai proposal, di lebih dari satu mesin, dan dengan waktu Proses untuk semua ukuran.
    \item \textbf{\emph{Nodata} sebelum inisialisasi.} Piksel \emph{nodata} sebaiknya dikeluarkan sebelum \emph{centroid} awal dipilih, agar $K$ besar tidak menghasilkan kelompok kosong.
    \item \textbf{Kanal tambahan dan validasi.} Kanal NIR (B08) dan SWIR dapat membantu memisahkan air dari bayangan serta permukiman dari tanah terbuka. Hasilnya perlu divalidasi terhadap peta acuan atau data lapangan sebelum dipakai sebagai klasifikasi.
\end{enumerate}
